
\subsubsection{6.1 Key Findings}

\textbf{The architecture-family adversarial fingerprint signal is large at base model scale, but statistically underpowered at N=9.} $\eta^{2}=0.293$ ($f^{2}\approx 0.41)$ is consistent across 83\% of attack categories. This is not a borderline effect: the estimate substantially exceeds Cohen's $f^2$>0.35 large-effect threshold. The finding is directionally consistent with Neerudu et al. (2023), who observed decoder robustness advantages over encoder models, and provides the first formal effect-size quantification of that phenomenon under multivariate conditions.

\textbf{adv\_rte (NLI binary entailment) produces the strongest family differentiation.} $\eta^{2}=0.592$ (p=0.075, near-significant at N=9) suggests that adversarial entailment detection concentrates architecture-family effects. This is consistent with the hypothesis that bidirectional and causal attention mechanisms diverge most on tasks requiring global whole-sentence semantic integration.

*\textit{The $\Delta$}-vector pipeline is operationally validated.*\textit{ All seven modules executed end-to-end, checkpoint recovery was exercised successfully, and the $9\times 6 \Delta$}-matrix was produced without errors.

\subsubsection{6.2 Limitations}

\textbf{Statistical underpowering (primary).} N=9 yields approximately 40\% power at $\eta^{2}=0.293$. The p=0.147 result is fully expected from the power analysis; it does not imply the effect is absent or small. The $\eta^{2}=0.293$ effect-size estimate stands as a practically meaningful finding independent of the p-value. The h-e1-v2 study with $N\geq 15$ directly addresses this limitation with approximately 80\% power.

\textbf{LOMO classification N-degeneracy.} LOMO=0.333 at exact chance. At N=3 per family, cosine k=1 nearest-neighbor classification is geometrically near-degenerate; at-chance LOMO is mathematically expected. The LOMO test at N<5 models per family is not a valid test of the classification hypothesis. Running LOMO at N=15 in h-e1-v2 is the direct test of whether accuracy rises above chance as degeneracy is resolved.

\textbf{Mechanism untested.} The attention-topology mediation hypothesis (planned as subsequent hypotheses h-m1 through h-m4 in the research program, covering $\Delta C$ extraction and attention-concentration analysis) was not executed in this experiment. The descriptive finding ($\eta^{2}=0.293)$ and methodological contribution (validated pipeline) are independent of mechanism confirmation.

\textbf{Encoder-decoder ANLI-R3 coverage gap.} T5-base and BART-base were fine-tuned only on SST-2, not MNLI. Their ANLI-R3 clean accuracy is near chance, producing $\Delta*\approx 0$ for both models and $\eta^{2}=0.000$ for the ANLI-R3 category. This is a protocol gap, not a genuine null result. The corrective step --- MNLI fine-tuning for T5 and BART --- is specified in h-e1-v2.

\textbf{CheckList partition skipped.} CheckList behavioral testing was not evaluated due to package unavailability. Cross-partition replication is therefore partial: results replicate across AdvGLUE (automatic attacks) and to a limited extent across ANLI-R3 (human-crafted, surrogate-free), but the CheckList behavioral partition is absent. Full cross-partition replication requires CheckList in h-e1-v2.

\textbf{AdvGLUE surrogate bias.} AdvGLUE attack examples were generated using encoder models (BERT/RoBERTa) as surrogate models, which may inflate encoder-family $\Delta$* values relative to decoder-only models. ANLI-R3 (surrogate-free, human-crafted) serves as a partial control, but enc\_dec coverage is missing from ANLI-R3 in this experiment. Surrogate-bias decomposition --- comparing $\eta^2$ across AdvGLUE, ANLI-R3, and CheckList once coverage is complete --- is specified as future work.

\textbf{Single-seed fine-tuning.} All models were evaluated with seed=42 only. Variance across seeds was not assessed in this proof-of-concept.

\textbf{Scale limitation.} Results are specific to base-scale models (~110--350M parameters). Whether architecture-family fingerprints persist at 7B+ scale, or under RLHF and instruction tuning, is an open empirical question not addressed here.

\subsubsection{6.3 Broader Impact}

The $\Delta$*-vector pipeline provides a practical tool for robustness auditing at the architecture-family level, potentially enabling more principled model selection for adversarially sensitive deployments. The validated seven-module codebase enables community replication and extension to new architectures, scales, and attack types.

Adversarial vulnerability characterization is dual-use. The characterization in this paper operates at the family level rather than the individual-model level, and all attack methods evaluated are publicly available benchmarks. No new attacks are introduced, and no novel individual-model vulnerabilities are disclosed beyond what the public benchmark literature already documents.

